\begin{table*}[t]
\centering\small
\caption{与学习对象相同的两模态梁：连续相位主频及其幅值依赖。}
\label{table:beam_kam}
\begin{tabular}{rrrr}
\toprule
$(r_1,r_2)$ & $(I_1^0,I_2^0)$ & $10^6(\omega_1-\nu_1)$ & $10^6(\omega_2-\nu_2)$ \\
\midrule
$(0.05,0.05)$ & $(0.00125,0.00125)$ & -2.0532 & -3.3235 \\
$(0.12,0.10)$ & $(0.00720,0.00500)$ & -8.9255 & -16.1932 \\
$(0.20,0.16)$ & $(0.02000,0.01280)$ & -23.3534 & -43.5138 \\
\bottomrule
\end{tabular}
\par\smallskip
\begin{minipage}{0.96\textwidth}\footnotesize
所有频率均对无量纲时间$\tau$计；物理角频率为$\omega_0\omega_j$。$I_j^0=r_j^2/2$仅标记谐振子参考初值，并非非线性系统的精确作用量。容差、采样及分窗检验采用$10^{-10}$的经验分辨阈值，该阈值不是严格误差界。
\end{minipage}
\end{table*}
